\documentclass[../../../include/open-logic-section]{subfiles}

\begin{document}

\olfileid{sth}{story}{rus}
\olsection{مفارقة راسل (مرة أخرى)}

عملنا في~\olref[sfr][][]{part} بنظرية ساذجة للمجموعات. لكن بحسب تصور
\emph{شديد} السذاجة، ليست المجموعات إلا امتدادات المحمولات. وتفرض هذه
الفكرة الساذجة المبدأ الآتي:

\begin{defish}
  \emph{الاستيعاب الساذج.} توجد~$\Setabs{x}{\phi(x)}$ لكل صيغة~$\phi$.
\end{defish}

هذا المبدأ مغرٍ، لكن يمكن إثبات عدم اتساقه. وقد رأينا ذلك في
\olref[sfr][set][rus]{sec}، غير أن النتيجة من الأهمية والوضوح بحيث
تستحق أن نكررها حرفيًا.

\begin{thm}[مفارقة راسل]
لا توجد مجموعة~$R = \Setabs{x}{x \notin x}$.
\end{thm}

\begin{proof}
إذا وجدت~$R = \Setabs{x}{x \notin x}$، فإن~$R \in R$ إذا وفقط إذا
كان~$R \notin R$، وهذا تناقض.
\end{proof}

اكتشف راسل هذه النتيجة في يونيو 1901. (لكنه لم يصغ المفارقة على الصورة
نفسها التي عرضناها للتو، لأنه كان ينظر في نظرية فرغه للمجموعات كما
عُرضت في~\emph{Grundgesetze}. وسنعود إلى ذلك في~\olref[blv]{sec}.)
وكتب راسل إلى فرغه في 16 يونيو 1902 موضحًا عدم الاتساق في نسقه.
للاطلاع على المراسلات وبعض الخلفية، انظر
\citet[pp.~124--8]{Heijenoort1967}.

ويجدر التشديد على أن هذا البرهان المؤلف من سطرين نتيجة من نتائج
\emph{المنطق المحض}. صحيح أننا استعملنا ضمنًا مسلّمة (غير منطقية؟)،
هي الامتدادية، في ترميزنا~$\Setabs{x}{x \notin x}$؛ إذ يُراد من
$\Setabs{x}{\phi(x)}$ أن تكون \emph{المجموعة الوحيدة} (بحسب
الامتدادية) للأشياء التي تحقق~$\phi$، إن وجدت. لكن يمكننا تجنب مجرد
الإيحاء بالامتدادية، بصياغة النتيجة هكذا:
\emph{لا توجد مجموعة تكون عناصرها بالضبط المجموعات التي لا تنتمي إلى
نفسها}. وليس لهذا صلة خاصة بالمجموعات. فكما لاحظ راسل نفسه، يؤدي
استدلال مماثل تمامًا إلى النتيجة: \emph{لا يوجد رجل يحلق بالضبط للرجال
الذين لا يحلقون لأنفسهم}. أو: \emph{لا يوجد كلب باغ يشم بالضبط كلاب
الباغ التي لا تشم نفسها}. وهكذا. وبصورة مخططية، ليست صورة النتيجة إلا:
\[
\lnot \exists x \forall z(Rxz \liff \lnot R zz).
\]
وهذه مجرد مبرهنة (مخططية) في منطق الرتبة الأولى. ولذلك لا يمكننا تفادي
مفارقة راسل بمجرد تعديل نظرية المجموعات؛ فهي تنشأ قبل أن \emph{نصل}
أصلًا إلى نظرية المجموعات. وإذا كنا سنستعمل منطق الرتبة الأولى
(الكلاسيكي)، فعلينا ببساطة أن \emph{نقبل} أنه لا توجد مجموعة
$R = \Setabs{x}{x\notin x}$.

والخلاصة هي الآتية. إذا أردنا قبول الاستيعاب الساذج مع
\emph{تفادي} عدم الاتساق، فلا يكفي أن نعدّل \emph{نظرية المجموعات}.
بل علينا أن نعيد بناء \emph{منطقنا}.

وقد قُدمت بالطبع نظريات للمجموعات ذات منطقيات غير كلاسيكية. لكنها،
على أقل تقدير، غير معيارية. والنهج المعياري إزاء مفارقة راسل هو معاملتها
بوصفها برهانًا مباشرًا على عدم الوجود، ثم محاولة تعلم كيفية التعايش
معها. وهذا هو النهج الذي سنتبعه.

\end{document}
